\chapter{Materiais, M\'etodos e Evolu\c{c}\~ao do Projeto}
\label{cap:metodologia}

\section{Natureza da pesquisa}
\label{sec:natureza}

Esta pesquisa \'e aplicada, predominantemente quantitativa e com forte componente tecnol\'ogico-experimental. \'E aplicada porque busca responder a um problema concreto de acompanhamento escolar. \'E quantitativa porque trabalha com dados hist\'oricos estruturados, m\'etricas preditivas e valida\c{c}\~ao temporal. Tamb\'em possui componente tecnol\'ogico porque n\~ao se limita \`a an\'alise conceitual: envolve a constru\c{c}\~ao efetiva de software, rotinas de extra\c{c}\~ao, pipeline de modelagem e relat\'orios operacionais.

\section{Fonte de dados e preparo inicial}
\label{sec:fonte-dados-preparo}

Os dados utilizados na vers\~ao publicada do projeto s\~ao dados escolares sintéticos para representar registros de alunos, notas, faltas, pagamentos e informa\c{c}\~oes acad\^emicas e administrativas sem expor pessoas reais. Esse conjunto foi estruturado para preservar o contrato p\'ublico da pipeline e a natureza temporal do problema, incluindo diferen\c{c}as de hist\'orico entre alunos e anos letivos. Como a avalia\c{c}\~ao temporal depende de conhecer a pr\'oxima nota observada, os experimentos finais utilizaram 2025 como ano de teste consolidado; dados de 2026 permanecem relevantes para continuidade operacional, mas nem sempre possuem alvo futuro completo para avalia\c{c}\~ao experimental. Essa decis\~ao ilustra um ponto importante do trabalho: nem todo dado dispon\'ivel na base publicada pode ser usado da mesma forma em uma avalia\c{c}\~ao preditiva honesta.

Os arquivos principais extra\'idos para a pipeline foram:
\begin{itemize}
    \item \texttt{aluno.csv};
    \item \texttt{media\_nota\_aluno.csv};
    \item \texttt{faltas\_aluno.csv};
    \item \texttt{pagamento\_aluno.csv};
    \item \texttt{responsaveis\_aluno.csv};
    \item \texttt{professor\_disciplina.csv}, quando o v\'inculo docente estava dispon\'ivel.
\end{itemize}

A descri\c{c}\~ao consolidada do contrato dessas entradas, da execu\c{c}\~ao experimental e dos artefatos gerados pela pipeline \'e apresentada na Se\c{c}\~ao~\ref{sec:config-experimental}.

\section{Estrat\'egia metodol\'ogica}
\label{sec:estrategia-metodologica}

O trabalho foi conduzido de forma incremental. Em vez de partir diretamente para uma solu\c{c}\~ao final, o projeto foi sendo refinado em ciclos de explora\c{c}\~ao, teste, diagn\'ostico de limita\c{c}\~oes e refatora\c{c}\~ao. Esse percurso n\~ao foi um desvio do objetivo do TCC; ao contr\'ario, foi justamente o processo que permitiu identificar quais escolhas se sustentavam diante do problema real e quais pareciam boas apenas em um plano inicial mais abstrato.

Para organizar a apresenta\c{c}\~ao desse percurso, tamb\'em foi adotada uma leitura inspirada no ciclo \textit{Cross-Industry Standard Process for Data Mining} (CRISP-DM) \cite{wirth2000}. Assim, o projeto pode ser compreendido em seis movimentos: entendimento do problema escolar, entendimento dos dados, prepara\c{c}\~ao da base, modelagem, avalia\c{c}\~ao e disponibiliza\c{c}\~ao dos resultados em relat\'orios e painel. Essa estrutura foi usada como material demonstrativo nos notebooks do reposit\'orio, sem substituir a arquitetura de produ\c{c}\~ao da aplica\c{c}\~ao. Em outras palavras, o CRISP-DM foi adotado aqui mais como lente de organiza\c{c}\~ao do percurso do que como descri\c{c}\~ao literal da implementa\c{c}\~ao final.

\section{Evolu\c{c}\~ao hist\'orica do projeto}
\label{sec:evolucao-historica}

\subsection{Fase inicial orientada a grafo}
\label{subsec:fase-grafo}

O projeto come\c{c}ou com uma hip\'otese baseada em grafos heterog\^eneos, motivada pela natureza relacional do dom\'inio escolar. Nessa etapa, foram produzidos m\'odulos para leitura de entidades, constru\c{c}\~ao de rela\c{c}\~oes e consultas em estruturas relacionais. A ideia era representar alunos, matr\'iculas, disciplinas, respons\'aveis, faltas e pagamentos como n\'os e arestas. Essa fase foi importante porque ajudou a explicitar a riqueza relacional do dom\'inio, mesmo que depois ela n\~ao tenha permanecido como base da solu\c{c}\~ao final.

\subsection{Fase explorat\'oria com testes cl\'assicos e PCA}
\label{subsec:fase-pca}

Na sequ\^encia, o projeto incorporou an\'alises estat\'isticas e testes com classificadores tradicionais, al\'em de redu\c{c}\~ao explorat\'oria de dimensionalidade com \textit{Principal Component Analysis} (PCA). Essa fase foi importante para entender a distribui\c{c}\~ao dos atributos e perceber que o problema exigia mais do que uma base est\'atica pivotada. Tamb\'em foi nesse momento que ficou mais claro que o desempenho do aluno n\~ao poderia ser tratado de forma satisfat\'oria sem considerar a ordem temporal das observa\c{c}\~oes.

\subsection{Fase temporal com LSTM}
\label{subsec:fase-lstm}

Como o problema possui natureza temporal, foi investigada tamb\'em uma abordagem sequencial com \textit{Long Short-Term Memory} (LSTM). Apesar de conceitualmente promissora, essa alternativa n\~ao se consolidou, sobretudo por depender de sequ\^encias mais homog\^eneas e longas do que as observadas na base. Na pr\'atica, a irregularidade dos hist\'oricos e a variedade de trajet\'orias entre alunos reduziram a ader\^encia dessa solu\c{c}\~ao ao escopo institucional do trabalho.

\subsection{Fase tabular com Random Forest}
\label{subsec:fase-random-forest}

Posteriormente, o projeto migrou para um desenho mais pragm\'atico com Random Forest. Essa fase representou um avan\c{c}o importante por aproximar a solu\c{c}\~ao de um problema tabular supervisionado. Ainda assim, permaneceram limita\c{c}\~oes relacionadas ao pivotamento excessivo, perda de temporalidade e uso incompleto de dados como faltas e pagamentos. Foi uma etapa decisiva porque mostrou que o projeto estava mais pr\'oximo de uma formula\c{c}\~ao tabular temporal do que de uma representa\c{c}\~ao fixa centrada apenas em atributos agregados.

\subsection{Refatora\c{c}\~ao para pipeline real}
\label{subsec:fase-pipeline-real}

O amadurecimento do trabalho levou primeiro \`a consolida\c{c}\~ao de uma pipeline intermedi\'aria de experimenta\c{c}\~ao, e depois \`a arquitetura final entregue no pacote \texttt{school\_predictor}. Essa arquitetura separou explicitamente:
\begin{itemize}
    \item a camada de acesso e conex\~ao com o banco;
    \item a prepara\c{c}\~ao local dos insumos e a gera\c{c}\~ao de dados sintéticos compat\'iveis com o contrato p\'ublico;
    \item a extra\c{c}\~ao dos dados para arquivos can\^onicos;
    \item a pipeline de modelagem e relat\'orios;
    \item a camada de visualiza\c{c}\~ao em dashboard.
\end{itemize}

Na vers\~ao final do reposit\'orio, a aplica\c{c}\~ao p\'ublica foi organizada em torno de \texttt{school\_predictor/}, enquanto detalhes locais de prepara\c{c}\~ao de insumos permaneceram fora do Git. Essa arquitetura final foi publicada no reposit\'orio oficial do projeto \cite{schoolpredictorgithub}, refor\c{c}ando a disponibilidade de uma vers\~ao acad\^emica consult\'avel e replic\'avel. Essa decis\~ao tornou a solu\c{c}\~ao mais adequada \`a publica\c{c}\~ao acad\^emica, concentrando a reprodu\c{c}\~ao no contrato dos arquivos e no uso de dados sintéticos. Mais do que uma escolha de organiza\c{c}\~ao de c\'odigo, isso foi parte do esfor\c{c}o para tornar o trabalho reproduz\'ivel sem depender de dados institucionais sens\'iveis.

\section{Crit\'erios de avalia\c{c}\~ao}
\label{sec:criterios-avaliacao}

Os crit\'erios de avalia\c{c}\~ao foram definidos em duas camadas. A primeira camada \'e estat\'istica e envolve MAE, RMSE, precision, recall e F1. A segunda camada \'e institucional e pergunta se a solu\c{c}\~ao consegue priorizar casos urgentes, gerar relat\'orios compreens\'iveis e apoiar diferentes atores escolares. Essa combina\c{c}\~ao foi necess\'aria porque um modelo matematicamente competente nem sempre \'e automaticamente \'util para a escola. Em um projeto como este, desempenho preditivo e utilidade operacional precisaram ser lidos em conjunto desde o in\'icio.
